% The oblique excess on an equidistant Cartesian mesh (two-dimensional cut, cell size h = 1).
% Interface: the circle with centre (0.818, 0.4) and radius 1.72, liquid inside. The plane of
% every interface cell is n.x = c with the exact normal at the midpoint of the cell's arc and c
% positioned by the cell's exact liquid fraction (computed on Lichtenberg; values below):
%   cell (0,1) alpha 0.99230  n (-0.42208, 0.90656) c 1.73636
%   cell (0,2) alpha 0.07211  n (-0.13452, 0.99091) c 1.98602
%   cell (1,1) alpha 0.91710  n ( 0.53689, 0.84365) c 2.48705   (N, the lateral neighbour)
%   cell (1,2) alpha 0.02960  n ( 0.23864, 0.97111) c 2.29799   (D1)
%   cell (2,0) alpha 0.51042  n ( 0.99815, 0.06075) c 2.53616
%   cell (2,1) alpha 0.16289  n ( 0.83382, 0.55204) c 2.60692   (D2; the donor D in panel a)
% The jump of N's and D2's planes on x = 2: z = 1.6752 (N) and 1.7015 (D2).
% Displacement per step: 0.56h along x (a), 0.4h per axis (b); the measured runs: 0.126h.
% Panels (c, d): the measured E_c of the studies sphereRotation3D_hexFineLine{Diag,X}{RingSwept,RDF},
% column 10 of staticFaceFractions.dat, copied from the cluster into data/oblique.
\begin{figure}[tb]
\centering
% the half-plane n.x <= c and the line n.x = c, for a clip
\newcommand{\obqHalf}[3]{({#3*#1-20*#2},{#3*#2+20*#1}) -- ({#3*#1+20*#2},{#3*#2-20*#1})
  -- ({(#3-20)*#1+20*#2},{(#3-20)*#2-20*#1}) -- ({(#3-20)*#1-20*#2},{(#3-20)*#2+20*#1}) -- cycle}
\newcommand{\obqLine}[3]{({#3*#1-20*#2},{#3*#2+20*#1}) -- ({#3*#1+20*#2},{#3*#2-20*#1})}
% the liquid inside the circle, clipped to the mesh
\newcommand{\obqLiquid}{%
  \begin{scope}
    \clip (0,0) rectangle (4,4);
    \fill[fluidcol] (0.818,0.4) circle[radius=1.72];
  \end{scope}}
% the interface and every cell's plane
\newcommand{\obqPlanes}{%
  \begin{scope}
    \clip (0,0) rectangle (4,4);
    \draw[ifacecol!55, line width=0.5pt] (0.818,0.4) circle[radius=1.72];
  \end{scope}
  \foreach \ci/\cj/\nx/\nz/\cc in {0/1/-0.42208/0.90656/1.73636, 0/2/-0.13452/0.99091/1.98602,
      1/1/0.53689/0.84365/2.48705, 1/2/0.23864/0.97111/2.29799,
      2/0/0.99815/0.06075/2.53616, 2/1/0.83382/0.55204/2.60692}{%
    \begin{scope}
      \clip (\ci,\cj) rectangle ++(1,1);
      \draw[ifacecol, line width=1.5pt] \obqLine{\nx}{\nz}{\cc};
    \end{scope}}}
% panel (b): the swept volumes of f1 and f2, their cuts and the faces
\newcommand{\obqCorner}{%
  \fill[sweptcol] (2,2) -- (2,3) -- (1.6,2.6) -- (1.6,1.6) -- cycle;
  \fill[sweptcol] (2,2) -- (3,2) -- (2.6,1.6) -- (1.6,1.6) -- cycle;
  \begin{scope}
    \clip (2,1.6) -- (2.6,1.6) -- (3,2) -- (2,2) -- cycle;
    \fill[gaincol] \obqHalf{0.83382}{0.55204}{2.60692};
  \end{scope}
  \begin{scope}
    \clip (1.6,1.6) -- (2,2) -- (1.6,2) -- cycle;
    \fill[mismcol] \obqHalf{0.23864}{0.97111}{2.29799};
    \fill[gaincol] \obqHalf{0.53689}{0.84365}{2.48705};
    \draw[ifacecol, dashed, line width=0.9pt] \obqLine{0.23864}{0.97111}{2.29799};
  \end{scope}
  \begin{scope}
    \clip (1.6,1.6) -- (2,1.6) -- (2,2) -- cycle;
    \fill[mismcol] \obqHalf{0.83382}{0.55204}{2.60692};
    \fill[gaincol] \obqHalf{0.53689}{0.84365}{2.48705};
    \draw[ifacecol, dashed, line width=0.9pt] \obqLine{0.83382}{0.55204}{2.60692};
  \end{scope}
  \begin{scope}
    \clip (1,1) rectangle (2,2);
    \draw[ifacecol, line width=1.5pt] \obqLine{0.53689}{0.84365}{2.48705};
  \end{scope}
  \begin{scope}
    \clip (2,1) rectangle (3,2);
    \draw[ifacecol, line width=1.5pt] \obqLine{0.83382}{0.55204}{2.60692};
  \end{scope}
  \draw[black!55, line width=0.5pt] (2,2) -- (2,3) -- (1.6,2.6) -- (1.6,1.6) -- cycle;
  \draw[black!55, line width=0.5pt] (2,2) -- (3,2) -- (2.6,1.6) -- (1.6,1.6);
  \draw[black!55, line width=0.4pt, densely dotted] (1.6,1.6) -- (2,2);
  \draw[facecol, line width=2.6pt] (2,2) -- (2,3);
  \draw[facecol, line width=2.6pt] (2,2) -- (3,2);}
\begin{tikzpicture}[scale=1.26, >=Stealth, font=\footnotesize]
  \colorlet{fluidcol}{blue!18}
  \colorlet{gaincol}{blue!45}
  \colorlet{ifacecol}{blue!65!black}
  \colorlet{facecol}{orange!80!black}
  \colorlet{bdrycol}{black!50}
  \colorlet{sweptcol}{black!12}
  \colorlet{mismcol}{red!60}
  % ------------------------------------------------ (a) translation along x
  \begin{scope}
    \obqLiquid
    \draw[bdrycol, line width=0.5pt] (0,0) grid[step=1] (4,4);
    % the swept volume of the face f = {x = 2, 1 <= z <= 2}: inside its donor D = (1,1)
    \fill[sweptcol] (1.44,1) rectangle (2,2);
    \begin{scope}
      \clip (1.44,1) rectangle (2,2);
      \fill[gaincol] \obqHalf{0.53689}{0.84365}{2.48705};
    \end{scope}
    \obqPlanes
    \draw[black!55, line width=0.5pt] (1.44,1) rectangle (2,2);
    \draw[facecol, line width=2.6pt] (2,1) -- (2,2);
    \node[facecol, anchor=west] at (2.04,1.88) {$f$};
    \node at (1.22,1.2) {$D$};
    \node at (2.8,1.2) {$R$};
    \draw[->, line width=0.8pt] (2.6,3.5) -- (3.16,3.5) node[midway, above=1pt] {$\mathbf{U}\Delta t$};
    \draw[<->, black!70, line width=0.5pt] (0,-0.18) -- (1,-0.18) node[midway, below=0pt] {$h$};
    \node[anchor=north] at (2,-0.42) {(a) along $x$};
  \end{scope}
  % ------------------------------------------------ (b) translation along the diagonal
  \begin{scope}[xshift=4.45cm]
    \obqLiquid
    \draw[bdrycol, line width=0.5pt] (0,0) grid[step=1] (4,4);
    \obqPlanes
    \obqCorner
    \draw[red!65!black, line width=0.6pt] (2,1.688) circle[radius=0.07];
    \node[facecol, anchor=west] at (2.04,2.62) {$f_1$};
    \node[facecol, anchor=south] at (2.62,2.03) {$f_2$};
    \node at (1.22,2.8) {$D_1$};
    \node at (2.8,1.2) {$D_2$};
    \node at (1.22,1.2) {$N$};
    \node at (2.8,2.8) {$R$};
    \draw[->, line width=0.8pt] (2.6,3.25) -- (3.0,3.65) node[pos=0.4, above left=0pt] {$\mathbf{U}\Delta t$};
    \node[anchor=north] at (2,-0.42) {(b) along the diagonal};
    % the square of N that moves into R, magnified three times to the right
    \draw[black!70, line width=0.4pt] (1.55,1.55) rectangle (2.05,2.05);
    \draw[black!40, line width=0.3pt] (2.05,2.05) -- (4.45,3.0) (2.05,1.55) -- (4.45,1.0);
    \begin{scope}[shift={(4.45,1.0)}, scale=4, shift={(-1.55,-1.55)}]
      \fill[white] (1.55,1.55) rectangle (2.05,2.05);
      \begin{scope}
        \clip (1.55,1.55) rectangle (2.05,2.05);
        \obqLiquid
        \draw[bdrycol, line width=0.5pt] (1,1) grid[step=1] (3,3);
        \obqCorner
        \draw[red!65!black, line width=0.6pt] (2,1.688) circle[radius=0.02];
      \end{scope}
      \draw[black!70, line width=0.4pt] (1.55,1.55) rectangle (2.05,2.05);
      \node[anchor=south] at (1.775,2.05) {$D_1$};
      \node[anchor=south] at (2.025,2.05) {$R$};
      \node[anchor=east] at (1.55,1.775) {$N$};
      \node[anchor=west] at (2.05,1.775) {$D_2$};
    \end{scope}
  \end{scope}
\end{tikzpicture}

\vspace{1mm}
\begin{tikzpicture}
  \begin{groupplot}[group style={group size=2 by 1, horizontal sep=1.4cm},
      width=6.4cm, height=3.7cm, xmin=0, xmax=0.45, ymin=0, ymax=0.065,
      xlabel={$t$}, ylabel={$E_c$}, tick label style={font=\scriptsize},
      label style={font=\footnotesize}, title style={font=\footnotesize, yshift=-1mm},
      xlabel style={yshift=1mm}, scaled y ticks=false, ytick={0,0.02,0.04,0.06},
      yticklabel style={/pgf/number format/fixed, /pgf/number format/precision=2},
      table/header=false, table/comment chars={\#}, table/x index=0, table/y index=9,
      legend style={font=\scriptsize, draw=none, fill=none, at={(0.02,0.98)}, anchor=north west},
      legend cell align=left]
    \nextgroupplot[title={(c) swept-face advection, ring normal}]
      \addplot[red!85!black, line width=1.2pt] table {data/oblique/lineDiagRingSwept_N64.dat};
      \addlegendentry{diagonal, $N = 64$}
      \addplot[red!85!black, line width=0.9pt, densely dashed] table {data/oblique/lineDiagRingSwept_N32.dat};
      \addlegendentry{diagonal, $N = 32$}
      \addplot[blue!80!black, line width=1.2pt] table {data/oblique/lineXRingSwept_N64.dat};
      \addlegendentry{$x$, $N = 64$}
      \addplot[blue!80!black, line width=0.9pt, densely dashed] table {data/oblique/lineXRingSwept_N32.dat};
      \addlegendentry{$x$, $N = 32$}
    \nextgroupplot[title={(d) plicRDF with isoAdvector}]
      \addplot[red!85!black, line width=1.2pt] table {data/oblique/lineDiagRDF_N64.dat};
      \addplot[red!85!black, line width=0.9pt, densely dashed] table {data/oblique/lineDiagRDF_N32.dat};
      \addplot[blue!80!black, line width=1.2pt] table {data/oblique/lineXRDF_N64.dat};
      \addplot[blue!80!black, line width=0.9pt, densely dashed] table {data/oblique/lineXRDF_N32.dat};
  \end{groupplot}
\end{tikzpicture}
\caption{The oblique excess on an equidistant Cartesian mesh of spacing $h$, two-dimensional cut. Every interface cell holds its plane (thick) with the exact normal, positioned by its exact volume fraction; orange: the faces whose fluxes are drawn; grey: their swept volumes ($\mathbf{U}\Delta t = 0.56h$, enlarged; the runs: $0.126h$); dark blue: the liquid crossing the face. (a) Along $x$ the swept volume of $f$ lies in its donor $D$ and is cut by $D$'s plane. (b) Along the diagonal (inset: the corner of $N$, magnified four times) the swept volumes of the faces $f_1$, $f_2$ of $R$ reach from their donors $D_1$, $D_2$ into the neighbour $N$; isoAdvector cuts the two triangles there by the donors' planes extended (dashed), the swept-face advection by $N$'s plane, and the cuts differ by the red regions; the red circle marks the jump between $N$'s and $D_2$'s planes. (c), (d) The measured $E_c$, \cref{eq:Ec}, of the translated sphere along the $x$--$z$ diagonal and along $x$, same initial state and time step.}
\label{fig:oblique}
\end{figure}
